\section{Введение}

Задача движения тела в гравитационном поле является одной из классических задач физики и математического моделирования. От точности ее решения зависит успех многих инженерных и научных проектов — от баллистических расчетов до планирования космических миссий. В простейшем случае, когда гравитационное поле считается однородным, а сопротивление воздуха отсутствует, движение описывается простыми кинематическими уравнениями. Однако в реальных условиях необходимо учитывать множество факторов: неоднородность гравитационного поля, сопротивление среды, вращение Земли и другие.

Целью данной работы является сравнительный анализ аналитических и численных методов решения задачи о движении тела в гравитационном поле. Для достижения поставленной цели необходимо решить следующие задачи:

\begin{enumerate}
    \item Вывести аналитические уравнения движения тела в однородном гравитационном поле;
    \item Разработать численный алгоритм (метод Эйлера и метод Рунге-Кутты 4-го порядка) для решения задачи в более сложной постановке;
    \item Реализовать алгоритмы на языке \texttt{Python};
    \item Провести сравнительный анализ полученных результатов;
    \item Визуализировать траектории движения при различных начальных условиях.
\end{enumerate}

Актуальность работы обусловлена широким применением методов математического моделирования в современной науке и технике \cite{bakhvalov2020, formalev2024}. Понимание преимуществ и недостатков различных численных схем позволяет выбирать оптимальный метод для каждой конкретной задачи.